\documentclass[12pt]{jlreq}
\usepackage{amsmath, amssymb}

\begin{document}

\newcommand{\C}{\mathbb{C}}
\newcommand{\R}{\mathbb{R}}
\newcommand{\Z}{\mathbb{Z}}
\newcommand{\Tr}{\operatorname{Tr}}
\newcommand{\Impart}{\operatorname{Im}}
\newcommand{\AF}{\operatorname{AF}}
\newcommand{\AX}{\operatorname{AX}}
\newcommand{\GL}{\operatorname{GL}}

正方行列 \(V\) に対して，\(V^*\) を \(V\) のエルミート共役とし，\(V\) の対角成分の和を
\(\Tr(V)\) と書いて，
\[
  \|V\|=\sqrt{\Tr(VV^*)}
\]
と定義する。
\(A\) を \(n\) 次の複素正方行列，\(A\) の特性方程式の解を重複を込めて
\(\lambda_1,\lambda_2,\ldots,\lambda_n\) とする。
\(A\) は，\(A=UTU^*\) の形に分解できる。ただし，\(U\) はユニタリ行列，\(T\) は上三角行列である。
\(T=D+M\) とし，\(D\) は対角行列で，\(M\) の対角成分はすべて \(0\) であるとする。
また，\(T^*T-TT^*\) の第 \((i,j)\) 成分を \(\tau_{ij}\) とする。このとき，次の (1)〜(4) に答えよ。

\begin{enumerate}
  \item \(\displaystyle \|M\|^2=\|A\|^2-\sum_{i=1}^{n}|\lambda_i|^2\) を示せ。
  \item \(\displaystyle \sum_{i=1}^{n}\tau_{ii}=0\) を示せ。
  \item 不等式
  \[
    \|M\|^2\le \tau_{22}+2\tau_{33}+3\tau_{44}+\cdots+(n-1)\tau_{nn}
  \]
  を示せ。
  \item 不等式
  \[
    \|A\|^2-\sum_{i=1}^{n}|\lambda_i|^2
      \le \sqrt{\frac{n^3-n}{12}}\ \|A^*A-AA^*\|
  \]
  を示せ。
\end{enumerate}

\end{document}
